\section*{Chapter \ref{ch:b3:x-ray-diffraction} --- Crystallography and X-ray Diffraction}

\begin{solution}{exo:b3:x-ray-diffraction:1}
$d = a/\sqrt N$: $d_{111} = \qty{325.66}{pm}$, $d_{200} = \qty{282.03}{pm}$, $d_{220} =
\qty{199.42}{pm}$.
\end{solution}

\begin{solution}{exo:b3:x-ray-diffraction:2}
FCC: the first reflections are (111) and (200). $\sin\theta = \lambda\sqrt N/2a$: $2\theta =
\ang{43.32}$ and \ang{50.45}.
\end{solution}

\begin{solution}{exo:b3:x-ray-diffraction:3}
P: all seven. I ($h + k + l$ even): 110, 200, 211, 220. F (unmixed): 111, 200, 220.
\end{solution}

\begin{solution}{exo:b3:x-ray-diffraction:4}
$\mathbf a^* = \mathbf e_x/a$, $\mathbf b^* = \mathbf e_y/b$: a rectangular lattice of sides $1/a$
and $1/b$, the long side of the direct lattice becoming the short one.
\end{solution}

\begin{solution}{exo:b3:x-ray-diffraction:5}
$\sin^2\theta$ ratios $1 : 2 : 3 : 4 : 5$, i.e. $N = 2, 4, 6, 8, 10$: body-centred (110,
200, 211, 220, 310). From the first line, $d_{110} = \lambda/2\sin(\ang{20.135}) =
\qty{223.77}{pm}$ and $a = d\sqrt2 = \qty{316.5}{pm}$: tungsten.
\end{solution}

\begin{solution}{exo:b3:x-ray-diffraction:6}
$Z = \rho N_Aa^3/M = 1.99 \times 6.022\times10^{23} \times (6.2929\times10^{-8})^3/74.6 = 4.0$.
\end{solution}

\begin{solution}{exo:b3:x-ray-diffraction:7}
$\beta = \qty{0.40}{\degree} = \qty{6.98e-3}{rad}$, $\cos\theta = \cos\ang{15.85} = 0.962$:
$L = 0.9 \times 0.15406/(6.98\times10^{-3} \times 0.962) = \qty{21}{nm}$.
\end{solution}

\begin{solution}{exo:b3:x-ray-diffraction:8}
$F_{hkl} = f_{\ce{Cs+}} + f_{\ce{Cl-}}(-1)^{h+k+l}$: $F_{100} = 54 - 18 = 36$, $F_{110} = 72$:
intensities in the ratio 1 : 4. The centre is occupied by a different ion from
the corners: the lattice is primitive (a body-centred lattice needs identical
contents at both points), and 100 is present.
\end{solution}

\begin{solution}{exo:b3:x-ray-diffraction:9}
Each atom has a copy shifted by $(\frac12,\frac12,0)$, multiplying $F$ by $1 +
\eu^{\iu\pi(h+k)} = 1 + (-1)^{h+k}$, zero for $h + k$ odd.
\end{solution}

\begin{solution}{exo:b3:x-ray-diffraction:10}
Quasicrystals have long-range order without periodicity: they have no
lattice, so the restriction, which concerns lattices, does not apply to them.
Their sharp diffraction spots led to defining a crystal as any solid with an
essentially discrete diffraction pattern, periodic or not.
\end{solution}

\begin{solution}{exo:b3:x-ray-diffraction:11}
The $c$ glide maps $(x, y, z)$ to $(x, \frac12 - y, \frac12 + z)$. For $k = 0$ the two
phase factors are $\eu^{2\pi\iu(hx + lz)}$ and $\eu^{2\pi\iu(hx + lz + l/2)} = (-1)^l
\eu^{2\pi\iu(hx+lz)}$: their sum vanishes for odd $l$.
\end{solution}

\begin{solution}{exo:b3:x-ray-diffraction:12}
An inversion, a mirror or a glide turns a molecule into its mirror image, so a
crystal containing one would contain both enantiomers. A pure enantiomer
crystallises only in the 65 \omterm{def:b3:x-ray-diffraction:space-group}{space groups} built from rotations, screw axes
and translations (Sohncke groups), such as \textit{P}$2_12_12_1$. By Friedel's law the
pattern of one enantiomer equals that of the other; the absolute
configuration is obtained from the small deviations caused by anomalous
scattering of heavier atoms.
\end{solution}

\begin{solution}{pb:b3:x-ray-diffraction:1}
\textbf{1.} $\theta = \ang{14.171}$, $d = \lambda/2\sin\theta = \qty{314.64}{pm}$.
\textbf{2.} 314.64, 222.49, 181.66, 157.32, 140.71, \qty{128.45}{pm}.
\textbf{3.} 1.000, 2.000, 3.000, 4.000, 5.000, 6.000.
\textbf{4.} 1, 2, 3, 4, 5, 6.
\textbf{5.} Yes, as (100), (110), (111), (200), (210), (211) of a primitive cubic cell
with $a' = \qty{314.6}{pm}$.
\textbf{6.} $N = 7$ (and 15, 23\dots), which is not a sum of three squares; the
first difference would come at the eighth line. Six lines cannot tell a P
cell of edge $a'$ from an F cell of edge $2a'$.
\textbf{7.} $N = 4, 8, 12, 16, 20, 24$: (200), (220), (222), (400), (420), (422).
\textbf{8.} $a = 2d_{200} = \qty{629.3}{pm}$.
\textbf{9.} \ce{KCl} (\qty{629.29}{pm}).
\textbf{10.} $F = 4[f(\ce{K+}) - f(\ce{Cl-})]$ for all-odd indices, and the two ions have 18
electrons each: the amplitudes cancel.
\textbf{11.} Yes for both: \ce{Na+} (10) and \ce{Cl-} (18), \ce{K+} (18) and \ce{Br-} (36)
scatter differently; \ce{NaCl}'s (111) would be at \ang{27.36}, \ce{KBr}'s at \ang{23.38},
both inside the recorded range.
\textbf{12.} X-rays see \ce{K+} and \ce{Cl-} as nearly identical; identical ions on a
rock-salt arrangement form a simple cubic array of edge $a/2$, an apparent cell.
\textbf{13.} Four \ce{KCl}.
\textbf{14.} $\rho = 4 \times 74.6/(6.022\times10^{23} \times (6.293\times10^{-8})^3) =
\qty{1.99}{g/cm^3}$.
\textbf{15.} Six and six (octahedra of the other ion).
\textbf{16.} $a/2 = \qty{314.6}{pm}$.
\textbf{17.} (440), $N = 32$, at \ang{87.65} (the all-odd (333)/(511) being absent).
\textbf{18.} Intensities change with preferred orientation of the crystallites,
absorption, the fall-off of scattering factors and thermal motion; positions
depend only on the lattice.
\textbf{19.} From $d = \lambda/2\sin\theta$, $\Delta d/d = -\cot\theta\,\Delta\theta$: the same angular error
gives a smaller relative error at large $\theta$.
\textbf{20.} $d_{420} = \lambda/2\sin\ang{33.190} = \qty{140.71}{pm}$, $a = d\sqrt{20} =
\qty{629.3}{pm}$.
\textbf{21.} $L = 0.9 \times 0.15406/(4.36\times10^{-3} \times \cos\ang{33.19}) = \qty{38}{nm}$.
\textbf{22.} No: $\beta = \ang{0.0095}$, far below a typical instrumental width.
\textbf{23.} \textbf{$a = \qty{629.3}{pm}$: the powder is potassium chloride}, its odd
reflections silenced by two ions with the same number of electrons.
\end{solution}
